%=============================================================================
%  chapter1.tex -- Reading the Light from a Divertor
%
%  Uses the macros defined in thesis_main.tex Part 2. Defines no new notation.
%  Every \section carries a % QUESTION: comment.
%
%  BEAT (writing_specification.md Sec.2): SETUP and COMPLICATION.
%  A spectroscopist inverts a line ratio against a table; the table assumes
%  equilibrium ionisation balance; a divertor is never in equilibrium. The
%  obvious suspect -- that the excited states cannot keep up -- is set up here
%  and overturned in Chapter 3. DO NOT RESOLVE THE REVERSAL IN THIS CHAPTER.
%
%  This chapter reports no number produced by this work.
%=============================================================================

\chapter{Reading the Light from a Divertor}
\label{ch:intro}

A physicist sits at a console and looks at a spectrum. An optical fibre viewing
the bottom of a fusion reactor has collected light along a narrow line of sight,
and a spectrometer has spread it out by wavelength. Two lines stand out: a red
one at \SI{656}{\nano\metre} and a blue-green one at \SI{486}{\nano\metre}, both
emitted by hydrogen atoms. She divides the brightness of the first by the
brightness of the second, looks the resulting number up in a table, and reads
off an electron temperature and an electron density. That is how one of the
least accessible regions of the machine gets measured. It is not the hottest
part; the core is far hotter. It is the part where the plasma meets a solid
surface, and the one whose conditions decide whether that surface survives.

This thesis is about the table. It was built by solving a model of hydrogen's
excited states to steady state, and it is indexed on two numbers: temperature
and density. A table can be indexed on two numbers only if every other quantity
the answer depends on has been eliminated by an assumption. The assumption in
this case is that ionisation and recombination in the plasma have come into
balance, not that they are nearly balanced, or changing slowly, but that the
plasma has \emph{finished changing}.

The plasma in question never finishes changing. It is the most violently
unsteady region of the machine. So the question is whether the table still
holds, and if not, by how much it is wrong and where.

That is all this chapter does: put the reader in front of the spectrum, show
where the table comes from, and expose the assumption buried inside it.
Section~\ref{sec:divertor} explains what the region being measured is for and
why a probe cannot be put into it. Section~\ref{sec:hydrogen_atom} builds the
hydrogen atom and its spectrum from scratch, because everything downstream is a
statement about two of its lines. Section~\ref{sec:spectrometer} says what the
instrument records and what it does not. Section~\ref{sec:inversion} follows the
chain from a measured ratio to a pair of plasma parameters.
Section~\ref{sec:the_assumption} is the hinge: it identifies the quantity the
table has eliminated and separates two approximations that are routinely
confused with one another. Section~\ref{sec:transients} asks whether the plasma
is ever in the state the table assumes. Section~\ref{sec:prior_work} says what
is already known and by whom, which is more than one might guess.
Section~\ref{sec:question} states the question this thesis answers and names the
suspect it went looking for.

%-----------------------------------------------------------------------------
\section{Why the divertor is hard}
\label{sec:divertor}
% QUESTION: what is a divertor for, and why must the plasma inside it be
% measured with light rather than with an instrument?

\emph{The question: what is the thing being measured, and why is measuring it
hard enough to be worth a thesis?}

Start with the machine. A tokamak is a doughnut-shaped vessel in which a gas of
hydrogen isotopes is heated until its electrons separate from its nuclei. The
resulting mixture of free electrons and bare nuclei is a plasma, and it is held
away from the walls by a magnetic field wound around the doughnut. If the plasma
gets hot enough and stays confined long enough, deuterium and tritium nuclei fuse
and release energy. ITER, under construction in southern France, is the machine
built to demonstrate this at reactor scale.

Confinement is never perfect, and this is the useful part. Particles and heat
leak steadily outward across the magnetic field, and once they cross the last
magnetic surface that closes on itself they are no longer trapped: they stream
along field lines instead. The layer outside that surface is called the
scrape-off layer (SOL), and in a modern tokamak the field lines in it are not
allowed to strike the wall at random. They are steered magnetically to a
dedicated set of water-cooled plates at the bottom of the vessel. Those plates,
together with the volume of plasma and gas in front of them, are the
\emph{divertor}.

A chemical engineer already owns the right picture for this. The divertor is the
quench-and-separation section bolted onto the outlet of a very hot reactor. It
receives the reactor's waste heat and its ash, the helium that fusion produces.
It cools the stream and pumps the ash away, and it must do so without being
destroyed by the duty.

Being destroyed by the duty is the entire difficulty. The exhaust power does not
spread itself over the plate; it arrives concentrated in a narrow strip where the
field lines land.

How concentrated, and what the target can take, are two different questions, and
the ITER divertor physics basis is careful to separate them in a way worth
repeating. The engineering standard is
\SI{10}{\mega\watt\per\square\metre}: it has long been asserted, in that
document's words, that this ``is a value which must be managed by the actively
cooled divertor technology'', and it is the qualification standard for
laboratory testing of the components \cite[Sec.~1]{Pitts2019}. \textbf{That
figure is not a prediction of what arrives.} The same section states plainly
that ``the value of target loading in ITER is not therefore a prediction of
SOLPS modelling, but rather the constraint imposed to restrict the acceptable
range of core plasma performance''. It is a budget, and the modelling is run so
as to stay inside it.

What arrives when the divertor is not held inside it is larger. At the lowest
divertor neutral pressure in the ITER SOLPS-4.3 database, about \SI{2}{\pascal},
the peak load on a shaped target is about
\SI{25}{\mega\watt\per\square\metre}, rising to about
\SI{40}{\mega\watt\per\square\metre} if the heat flux channel turns out to be as
narrow as \SI{1}{\milli\metre} \cite[Sec.~4]{Pitts2019}. Against those, a
criterion based on avoiding tungsten recrystallisation over the roughly $2000$
hours of high-performance exposure expected of the first divertor gives
$16\pm\SI{2.5}{\mega\watt\per\square\metre}$, and
critical-heat-flux tests scaled to the final monoblock design put the ceiling on
the top surface at about \SI{20}{\mega\watt\per\square\metre}
\cite[Sec.~4]{Pitts2019}. The distance between what arrives and what the surface
tolerates is the reason a divertor is a research problem and not a plumbing
one, and closing it means removing the power before it reaches the surface at
all.

The way out is not a better plate material. It is to stop the plasma before it
arrives. If enough neutral hydrogen gas is allowed to build up in front of the
target, the plasma ions collide with it: in a charge-exchange collision an ion
takes an electron from a neutral atom and becomes a fast neutral itself, carrying
its momentum out of the plasma stream, and the energy that remains is radiated
away by excited atoms rather than deposited on the plate. Over the last few
centimetres the plasma pressure and the heat flux both collapse. The plasma is
then said to be \emph{detached} from the target: the region where ionisation
happens lifts off the plate and sits in the volume in front of it. Detachment is
not a fault to be avoided. It is the intended operating regime of a reactor
divertor.

Detachment is also what makes the atomic physics of hydrogen unavoidable. A
detached divertor plasma is cold, of order a few electronvolts, where one
electronvolt corresponds to about \SI{11600}{\kelvin}. At a few electronvolts
hydrogen is only partly ionised. Neutral atoms and ions coexist in
comparable numbers, recombination competes with ionisation on equal terms, and
the light emitted by neutral hydrogen becomes the brightest thing in the visible
spectrum. The plasma one most needs to characterise is therefore exactly the
plasma whose light carries the most information.

What conditions are those, in numbers? Edge modelling of the ITER baseline
burning plasma gives target electron temperatures spanning roughly one to a few
tens of electronvolts across the design database, with the large majority of
cases below \SI{15}{\electronvolt} \cite{Stangeby2023}, and
falling below \SI{1}{\electronvolt} across half a metre of each target in the
most strongly detached solutions obtained so far \cite{Lore2022}. The near-target
electron density under partial detachment is expected to reach about
\SI{e15}{\per\cubic\centi\metre} \cite{Krasheninnikov2017}. The same work
makes the point this thesis depends on, that at such densities radiative decay is
no longer fast compared with electron-impact processes, so the excited states
must be treated with a full collisional--radiative calculation rather than a
simple equilibrium formula.

Every calculation in this thesis is performed on a fixed grid spanning
$\Te = 1$--\SI{10}{\electronvolt} in electron temperature and
$\nel = 10^{12}$--\SI{e15}{\per\cubic\centi\metre} in electron density. The
upper density matches the expected near-target value cited above and the
temperature range covers the cold, detached end of the modelled database, which
is where the atomic physics of neutral hydrogen governs the emission. It does not
cover the hotter attached cases, and it is a modelling choice rather than a
prediction: it should be read as the range over which the model was exercised,
not as a claim about what any particular machine will do.
The lower end deserves a sharper statement than that, because it is below
anything the sources read for this chapter describe. Across the $53$ ITER
SOLPS-4.3 cases spanning a wide range of fuel and impurity throughput, the
upstream separatrix density is
\SIrange{3e19}{6e19}{\per\cubic\metre} \cite[Sec.~3.1.1]{Pitts2019}, that is
\SIrange{3e13}{6e13}{\per\cubic\centi\metre}, and a divertor density is higher
than an upstream one rather than lower. \textbf{The bottom of this grid,
\SI{e12}{\per\cubic\centi\metre}, is therefore an extension below the modelled
range and is not defended as an ITER divertor condition.} It is kept because the
coronal limit is approached there and the atomic physics is at its cleanest,
which makes it useful for testing the model even where it is not useful for
predicting a machine. Results at that edge are labelled throughout.

It would be simpler to put an instrument into the plasma, and ITER will do
exactly that. A Langmuir probe is a small electrode exposed to the plasma and
biased to a range of voltages; the current it draws as a function of bias gives
the local electron temperature and density. ITER carries a designed divertor
Langmuir probe system whose stated purpose is to measure plasma parameters at the
target plates, and whose thermomechanical analysis concludes that the sensors
survive the environment \cite{Zhao2020}. The claim that probes cannot survive in
a reactor divertor is often made and is not correct.

The limitation is not survival but what a probe can see. It is mounted flush in
the target surface, so it reports conditions in the thin charged layer at the
plate rather than in the volume in front of it, and the volume in front of it
is precisely where a detached plasma does its recombining. It perturbs what it
measures. Its interpretation relies on sheath theory that is least trustworthy in
exactly the cold, strongly recycling conditions of interest. And the coverage is
necessarily sparse: a set of sensors embedded in a handful of divertor cassettes
samples points on a surface, not a volume.

Light has neither limitation, and one decisive advantage: it leaves the plasma by
itself. A window, a fibre, a spectrometer and a detector make a measurement that
touches nothing, perturbs nothing, and samples the emitting volume rather than
the surface. Spectroscopy is therefore complementary to probe measurement rather
than a replacement for it, and it is the only one of the two that reports on the
recombining volume where a detached plasma does its interesting work. This is why
hydrogen line emission is a workhorse divertor diagnostic, and why the chain of
reasoning that converts photons into $\Te$ and $\nel$ carries so much weight. The
next section builds the atom that emits them.

%-----------------------------------------------------------------------------
\section{Light from a hydrogen atom}
\label{sec:hydrogen_atom}
% QUESTION: what is a spectral line, and which two lines is this thesis about?

\emph{The question: where do the two lines in the opening paragraph come from,
and what would make their brightnesses differ?}

A hydrogen atom is one electron bound to one proton. The electron cannot take
any energy it likes: the bound states form a discrete ladder, labelled by an
integer $n = 1, 2, 3, \dots$ called the principal quantum number, with energies
\begin{equation}
  E_n \;=\; -\,\frac{\SI{13.6}{\electronvolt}}{n^{2}} .
  \label{eq:bohr_levels}
\end{equation}

\noindent Here $E_n$ is the energy of the level in electronvolts, measured from
the free-electron state at $E = 0$, so all bound levels are negative and $n = 1$,
the ground state, is the deepest at \SI{-13.6}{\electronvolt}. What
Eq.~\eqref{eq:bohr_levels} does physically is fix the whole spectrum of hydrogen
with one constant: the levels crowd together as $n$ grows, and \SI{13.6}{\electronvolt}
is the energy needed to remove the electron entirely from the ground state, which
is called the ionisation energy.

Each level $n$ is not a single state. The electron's orbital angular momentum is
also quantised, by an integer $\ell = 0, 1, \dots, n-1$, and states of different
$\ell$ have historically been labelled by letters: $\ell = 0$ is written S,
$\ell = 1$ is P, $\ell = 2$ is D, $\ell = 3$ is F. So the $n = 2$ level contains
a $2s$ state and a $2p$ state, the $n=3$ level contains $3s$, $3p$ and $3d$, and
so on. In an isolated hydrogen atom all the $\ell$ states of a given $n$ have
almost exactly the same energy, which is why they can often be lumped together,
and Chapter~\ref{ch:atoms} explains the one place in this problem where
lumping them destroys the physics. The number of independent quantum states in a
level is its statistical weight $g$, equal to $2n^{2}$ for hydrogen when electron
spin is counted; a level with a larger $g$ has more ways of being occupied and is
correspondingly favoured in thermal equilibrium.

An atom in an excited level does not stay there. It can fall to a lower level and
emit the energy difference as a single photon. This is spontaneous emission, and
it is a first-order rate process exactly like a unimolecular decomposition: the
probability per unit time that an atom in level $q$ decays to a lower level $p$
is a constant $A_{qp}$, in \si{\per\second}, called the Einstein coefficient for
spontaneous emission. The number of such decays per unit volume per unit time is
therefore $A_{qp}n_q$, where $n_q$ is the population of level $q$ in
\si{\per\cubic\centi\metre}. Not every pair of levels is connected: quantum
mechanics permits a strong (electric-dipole) transition only when $\ell$ changes
by exactly one, a restriction called a selection rule, and Chapter~\ref{ch:atoms}
shows that the levels which \emph{cannot} decay this way are the ones that
dominate this problem.

Because the energy released is fixed by Eq.~\eqref{eq:bohr_levels}, every
transition emits at one wavelength, and the emission appears in a spectrograph as
a narrow \emph{spectral line}. Its brightness is proportional to the population
of the upper level. That single sentence is why spectroscopy is a diagnostic at
all: the light tells you how the atoms are distributed over their levels, and the
distribution is set by the plasma conditions.

The transitions that end on $n = 2$ fall in the visible, which is convenient for
windows, fibres and detectors, and they are collectively called the Balmer
series. The two brightest are the subject of this thesis:
\begin{align}
  H_\alpha:\quad n=3 \to 2, \qquad
  \Delta E &= \SI{13.6}{\electronvolt}\left(\tfrac{1}{4}-\tfrac{1}{9}\right)
            = \SI{1.89}{\electronvolt} \;\Rightarrow\; \SI{656}{\nano\metre},
  \label{eq:halpha}\\
  H_\beta:\quad n=4 \to 2, \qquad
  \Delta E &= \SI{13.6}{\electronvolt}\left(\tfrac{1}{4}-\tfrac{1}{16}\right)
            = \SI{2.55}{\electronvolt} \;\Rightarrow\; \SI{486}{\nano\metre}.
  \label{eq:hbeta}
\end{align}

\noindent The arithmetic is worth doing once by hand, because it shows that the
red line and the blue-green line of the opening paragraph are not empirical
facts to be looked up: they follow from Eq.~\eqref{eq:bohr_levels} and the
relation $\lambda = hc/\Delta E$, with $hc = \SI{1240}{\electronvolt\nano\metre}$.

What these two lines have in common is their lower level. Both end on $n = 2$,
so the ratio of their brightnesses does not depend on anything about $n = 2$: it
depends only on how the atoms are shared between the $n = 3$ and $n = 4$ levels.
The whole of this thesis is a statement about that one ratio, and
Section~\ref{sec:spectrometer} explains why it, and not the brightness of either
line by itself, is what a divertor spectroscopist can actually use.

%-----------------------------------------------------------------------------
\section{What a spectrometer records}
\label{sec:spectrometer}
% QUESTION: what does the instrument actually give you, and what does it not?

\emph{The question: given that a spectrometer measures brightness, what can and
cannot be extracted from that measurement?}

Picture the geometry first. The instrument looks along a straight line through
the plasma and collects the light emitted anywhere along it, from a small solid
angle. It cannot tell where along the line a photon came from. Whatever it
reports is a sum over the whole chord.

Formally, if each decay from level $q$ to level $p$ releases a photon of energy
$h\nu_{qp}$ in joules and the decays occur at $A_{qp}n_q$ per unit volume per
unit time, the detector records the brightness
\begin{equation}
  I_{qp} \;=\; \frac{h\nu_{qp}}{4\pi}\int_{\text{chord}} A_{qp}\,n_q(s)\,\mathrm{d}s ,
  \label{eq:brightness}
\end{equation}

\noindent in \si{\watt\per\square\metre\per\steradian}, where $s$ is distance
along the line of sight in \si{\centi\metre} and the factor $4\pi$ accounts for
the fact that an atom radiates into all directions equally, so only that fraction
of the emission reaches the detector. Equation~\eqref{eq:brightness} says the
brightness is the emission rate per unit volume, integrated over the length of
plasma the instrument happens to be looking through.

It also assumes that a photon, once emitted, escapes without being reabsorbed by
another atom on its way out. A plasma for which that is true is called
\emph{optically thin} at that wavelength. For the Balmer lines in a divertor this
is usually a fair assumption, because the lower level $n=2$ holds very few atoms
to absorb them; for transitions that end on the ground state it is often badly
false, since the ground state holds almost all of them. Chapter~\ref{ch:limits}
takes that up, quantifies it for this model, and it is the single most important
restriction on what follows.

Now the discouraging part. Equation~\eqref{eq:brightness} has one measured number
on the left and, on the right, the population we want multiplied by a path length
nobody knows, integrated over an emitting region whose extent nobody knows
either. Worse, $n_q$ is itself proportional to the amount of neutral hydrogen in
the viewing volume, and that is set by how much gas recycles off the target and
how it moves, neither of which the spectrometer sees. An absolutely calibrated
brightness of a single line therefore determines a product of unknowns and none
of them individually.

Take the ratio of two lines and most of this cancels. Define the observable used
throughout this thesis:
\begin{equation}
  \Robs \;\equiv\; \frac{I_{H\alpha}}{I_{H\beta}}
        \;=\; \frac{\nu_{32}\,A_{32}}{\nu_{42}\,A_{42}}
              \cdot\frac{\int n_3\,\mathrm{d}s}{\int n_4\,\mathrm{d}s},
  \label{eq:ratio_def}
\end{equation}

\noindent where $n_3$ and $n_4$ are the populations of the $n=3$ and $n=4$
levels in \si{\per\cubic\centi\metre} and $\Robs$ is dimensionless. The
prefactor is a ratio of atomic constants, known to high accuracy. What
Eq.~\eqref{eq:ratio_def} does is remove the geometry: provided the two lines are
emitted by the same volume, the path length divides out, the absolute neutral
density divides out, and the absolute calibration of the instrument divides out
with them. That proviso is the substantive assumption here, and a good one for
two lines of the same series from the same species. Only the \emph{relative}
sensitivity of the detector at two nearby visible wavelengths is needed, which is
a far easier measurement to make and to maintain.

What survives is a pure ratio of two excited-state populations. That is the whole
appeal of a line ratio, and it is why $\Robs$, rather than any single brightness,
is the quantity a divertor spectroscopist works with. It is also where the burden
of interpretation now sits. A line ratio contains no geometry and no absolute
density, so every piece of physics in it is compressed into one question:
\emph{what sets the ratio of the $n=3$ and $n=4$ populations?} The next section
answers that with a model, and the section after it examines what the model had
to assume.

%-----------------------------------------------------------------------------
\section{From a ratio to a temperature}
\label{sec:inversion}
% QUESTION: how is a measured ratio turned into a temperature and a density, and
% where does the model enter that chain?

\emph{The question: what stands between the number the spectrometer produces and
the $(\Te,\nel)$ the physicist writes down?}

An atom in a plasma is pulled in several directions at once. Free electrons
collide with it and knock its bound electron up the ladder of
Eq.~\eqref{eq:bohr_levels}, or right off it, which is \emph{ionisation}. Ions
capture free electrons back into bound levels, which is \emph{recombination};
this happens either with the excess energy carried off by a photon (radiative
recombination) or with it carried off by a second electron that happens to be
present (three-body recombination, which therefore needs two electrons and
becomes important at high density). Meanwhile every excited atom is decaying
radiatively on its own clock, and every decay from a high level cascades down
through the intermediate ones. No single one of these processes dominates in a
divertor. All of them run at once, and their rate coefficients span more than
ten orders of magnitude, so which of them matters depends on the level and on
the conditions rather than being settled once.

The object that keeps track of all of it is a \emph{collisional--radiative} (CR)
model. It is nothing more exotic than a population balance written for every
level of the atom: accumulation equals everything arriving minus everything
leaving, with one term per atomic process and a measured or computed rate
coefficient attached to each. The framework was set out for hydrogenic ions by
Bates, Kingston and McWhirter in 1962 \cite{Bates1962}, whose central point was
that neither of the two classical limits describes a laboratory plasma: not
coronal equilibrium, where the density is so low that every collisional
excitation is followed by a radiative decay before anything else happens, and not
local thermodynamic equilibrium (LTE), where collisions are so frequent that the
level populations take their thermal equilibrium values. Real plasmas sit between
the two, and the populations there have to be computed.

A chemical engineer will recognise this object on sight. It is a reaction network
in a well-mixed vessel: one species per energy level, one reaction per atomic
process, a feed stream, and a set of coupled first-order rate equations with a
rate-constant matrix. Chapter~\ref{ch:theory} builds it explicitly and pushes
that correspondence considerably further than a metaphor would go.

Solve the CR model at a fixed $\Te$ and $\nel$ and it returns the population of
every level, hence $n_3$, hence $n_4$, hence $\Robs$ through
Eq.~\eqref{eq:ratio_def}. Do that on a grid of $(\Te,\nel)$ and the result is a
table, or equivalently a family of contours in the $(\Te,\nel)$ plane along which
$\Robs$ is constant. Inverting a measurement means locating the measured value of
$\Robs$ on that map. One ratio is one constraint on two unknowns, so practice
supplies a second: another line ratio, the Stark broadening of a line by the
plasma's own microscopic electric fields, continuum emission, or an independent
estimate of one of the two parameters. The intersection of two contours fixes the
point. Figure~\ref{fig:diagnostic_chain} draws the chain end to end, from the
recorded spectrum to the reported $(\Te,\nel)$, and marks the step at which the
table's own assumption enters.

\begin{figure}[tbp]
  \centering
  \includegraphics[width=\textwidth]{fig1_1_diagnostic_chain}
  \caption{\CapFigDiagnosticChain}
  \label{fig:diagnostic_chain}
\end{figure}

This is not a hypothetical procedure. Balmer-spectrum analysis against a CR
model has been used to infer divertor conditions since at least the Alcator
C-Mod measurements of Pigarov, Terry and Lipschultz, who matched calculated and
measured Balmer spectra, with populations of the excited states supplied by a
collisional--radiative model, to determine the plasma parameters in the
recombining region of a divertor at electron temperatures near one electronvolt
\cite{Pigarov1998}. The modern versions of the technique are more elaborate but
the chain is the same one.

In practice the table is rarely built by hand. The Atomic Data and Analysis
Structure (ADAS) distributes precomputed coefficients for exactly this purpose,
tabulated over $(\Te,\nel)$ \cite{Summers2006, ADAS}. What matters for the
argument of this thesis is a detail of how ADAS stores them: for each line it
supplies \emph{two} emissivity coefficients, one for the emission driven by
excitation out of the neutral ground state and one for the emission driven by
recombination from the ion. The database keeps the two apart deliberately. The
next section is about what happens when they are combined.

%-----------------------------------------------------------------------------
\section{The assumption inside the table}
\label{sec:the_assumption}
% QUESTION: what must a table indexed on (Te, ne) alone assume, and are there
% really two separate assumptions rather than one?

\emph{The question: a table has two axes; how many quantities does the answer
actually depend on?}

Take the physical picture before any algebra. An atom sitting in the $n=3$ level
got there by one of two routes. Either it started as a neutral atom in the ground
state and an electron collision kicked it up, fed \emph{from below}, or it
started as an ion and captured a free electron into a high level, from which it
cascaded down, fed \emph{from above}. Those are two different supply lines
drawing on two different reservoirs: the population of ground-state neutral
atoms, $\ngs$, and the population of ions, $\nion$. Every excited level is
filled by some mixture of the two.

Because the population balances are linear in the populations, the mixture is
exact rather than approximate. For every level $p$,
\begin{equation}
  n_p \;=\; \underbrace{a_p\,\ngs}_{\text{fed from below}}
        \;+\; \underbrace{c_p\,\nion}_{\text{fed from above}},
  \label{eq:two_supplies_intro}
\end{equation}

\noindent where $n_p$, $\ngs$ and $\nion$ are in \si{\per\cubic\centi\metre},
the dimensionless coefficient $a_p$ is the number of atoms in level $p$ supported
per ground-state atom, and $c_p$ is the number supported per ion. What
Eq.~\eqref{eq:two_supplies_intro} does is separate the atomic physics from the
plasma composition: $a_p$ and $c_p$ depend on $\Te$ and $\nel$ and on nothing
else, while $\ngs$ and $\nion$ describe how much of each reservoir is present.
This split is as old as the CR model itself. It is the population-coefficient
decomposition of Bates, Kingston and McWhirter, written by Fujimoto and McWhirter
in the normalised form Chapter~\ref{ch:theory} adopts
\cite{Bates1962, FujimotoMcWhirter1990, Fujimoto2004}. Its two limits have names
that will recur: a plasma in which the first term dominates is called
\emph{ionising}, one in which the second dominates is called \emph{recombining},
and a divertor moves between the two as its conditions change. The crossover is
not identified with any named operating regime here; Chapter~\ref{ch:limits}
gives the densities and explains why that identification does not hold.

Now put Eq.~\eqref{eq:two_supplies_intro} into the observable. For the two Balmer
levels,
\begin{equation}
  \Robs \;\propto\; \frac{a_3\,\ngs + c_3\,\nion}{a_4\,\ngs + c_4\,\nion}
         \;=\; \frac{a_3\,(\ngs/\nion) + c_3}{a_4\,(\ngs/\nion) + c_4}.
  \label{eq:ratio_three_args}
\end{equation}

\noindent The second form is the important one: dividing through by $\nion$ shows
that the observable depends on the atomic coefficients, which are functions of
$(\Te,\nel)$, \emph{and} on the ratio $\ngs/\nion$, which is not. That ratio is a
third, independent quantity. It says how much neutral hydrogen is present per
ion, and it is fixed by the plasma's recent history and by transport, not by the
local temperature and density.

A table with two axes must therefore dispose of the third quantity, and there are
only two honest ways to do it. One is to work in a region of the operating range
where a single supply dominates \emph{both} levels, so that the third argument
cancels between numerator and denominator of Eq.~\eqref{eq:ratio_three_args} and
the table is right whatever the neutral fraction does. The other is to adopt a
closure, a rule that fixes $\ngs/\nion$ in terms of $\Te$ and $\nel$. The rule
used in practice is \emph{equilibrium ionisation balance}: solve the CR system
all the way to steady state, and take the neutral-to-ion ratio it predicts for
those local conditions. The state it predicts is called collisional--radiative
equilibrium (CRE), and it is worth being clear that this is a much weaker
condition than full thermodynamic equilibrium: it means only that the CR rate
equations have stopped changing, not that the populations are thermal.

That closure is an assumption about the plasma, not a piece of atomic physics.

It is important to say who is and is not being criticised, because the databases
are not at fault. ADAS keeps the excitation-driven and recombination-driven
coefficients separate precisely so that a user can weight them by the $\ngs$ and
$\nion$ appropriate to the plasma in hand \cite{Summers2006}. The most careful
divertor analyses do exactly that, and treat the neutral content as an unknown
rather than a given: Verhaegh and co-workers, analysing Balmer emission from the
divertor of the TCV tokamak, take the neutral fraction $\ngs/\nel$ to be a
quantity that is estimated rather than measured, assign it a broad prior spanning
more than an order of magnitude, and propagate that uncertainty through the whole
inference by Monte Carlo so that every inferred quantity emerges as a probability
distribution rather than a number \cite{Verhaegh2019}. In that community, the
dependence written in Eq.~\eqref{eq:ratio_three_args} is known and is handled,
and this thesis does not discover it.

The same paper is worth one more sentence, because it draws a boundary this
thesis has to respect. Verhaegh and co-workers restrict their analysis to the
higher Balmer lines, $n \geq 5$, on the grounds that the lower ones, $H_\alpha$
above all, are contaminated by emission from molecular processes that their
atomic treatment does not model \cite{Verhaegh2019}. This thesis studies
$H_\alpha$ and $H_\beta$, which is to say precisely the two lines that careful
practice avoids in detached conditions, and it does so with a model containing no
molecules. Section~\ref{sec:prior_work} returns to this and
Chapter~\ref{ch:limits} quantifies what it costs; it is stated here so that no
reader reaches those chapters feeling it was kept back. The assumption enters at
the moment a measured ratio is inverted against a two-parameter lookup, because a
two-parameter lookup returns exactly one value of $\Robs$ per $(\Te,\nel)$ pair
and therefore requires the two supplies to have been combined in some fixed
proportion.

There are now two approximations in play, and keeping them apart is the single
most important thing this thesis does.

\begin{description}
\item[The quasi-steady-state (QSS) approximation.] That the \emph{excited} levels
  have already reached steady state with whatever reservoirs are present at that
  instant. This is the standard shortcut of CR theory, and it is what makes
  $a_p$ and $c_p$ in Eq.~\eqref{eq:two_supplies_intro} functions of $(\Te,\nel)$
  alone. \textbf{The approximation being tested in this thesis is the Bodenstein
  steady-state approximation, applied to atomic level populations instead of
  radical concentrations.} It is the same approximation the reader has already
  met for reactive intermediates, with the same justification, that the
  intermediates turn over far faster than anything else in the system. In
  principle it has the same failure mode too.

\item[The ionisation-balance assumption.] That the \emph{reservoirs themselves},
  $\ngs$ and $\nion$, stand in the proportion that CRE would give for the local
  $\Te$ and $\nel$. This is what removes the third argument from
  Eq.~\eqref{eq:ratio_three_args} and allows the table to have two axes.
\end{description}

These are statements about different populations relaxing at different speeds,
and only the first is what the literature normally means by ``the validity of
QSS''. A two-parameter table needs both to hold. Chapter~\ref{ch:theory} makes
the distinction exact, defines a separate error measure for each, $\epsQSS$ and
$\epsCRE$, and shows why conflating them makes it possible to measure the second
and report it as the first. Whether either holds in a real divertor is the
subject of the next section.

%-----------------------------------------------------------------------------
\section{A plasma that never finishes changing}
\label{sec:transients}
% QUESTION: is a divertor plasma ever in the equilibrium state the table
% assumes?

\emph{The question: how long does the plasma get between disturbances, and is
that long enough?}

The ionisation-balance assumption would be harmless if divertor plasmas sat
still. They do not, and they are not meant to.

The largest regular disturbance is the edge-localised mode (ELM). The confined
plasma maintains a steep pressure gradient at its edge; periodically that
gradient becomes unstable, the edge relaxes, and a burst of particles and energy
is expelled into the SOL and flushed down the field lines onto the divertor.

The numbers that matter here are the ones describing a single event, and they are
short. In present-day machines an ELM heat pulse at the target lasts between
about \SI{0.75}{\milli\second} and \SI{3}{\milli\second}, and the scaling to
ITER gives a deposition rise time of \SI{250}{\micro\second} with a decay
roughly twice that, the pulse being well approximated by a triangular waveform;
each event carries away between two and ten per cent of the energy stored in the
plasma edge \cite{Eich2017}. So the conditions in the divertor are reset over a
few hundred microseconds, repeatedly.

How often they are reset is a more delicate question, and the usual figure is
weaker than it looks. The expected natural
ELM frequency in ITER is of order \SI{1}{\hertz} at the full plasma current of
\SI{15}{\mega\ampere}, rising to a few hertz at reduced current and scaling
roughly as the inverse square of the current \cite{Kirk2013}; present-day
machines, being smaller, run at tens to hundreds of hertz. The frequently quoted
range of one to two hertz for ITER traces, as far as could be established for
this chapter, only to a conference paper that was not peer reviewed, and it is
not repeated here for that reason. \todo{If a refereed derivation of the ITER
natural ELM frequency exists, cite it here in place of Kirk \emph{et al.} (2013),
whose statement of the number sits in an introduction rather than in a result.}

The distinction matters for the argument. An ELM is not a gentle modulation that
the plasma tracks quasi-statically; it is a step change delivered in a few
hundred microseconds, after which the plasma is left to recover. The relevant
question is therefore not how often the steps come but how long the recovery
takes relative to a single step, and that is a question about the plasma's own
internal clocks.

An ELM is not the only transient. The detachment front itself moves: as the
neutral pressure in the divertor is raised, the ionisation region lifts off the
plate and travels back up the field lines, and where it settles is controlled by
feedback on a gas valve. Operating near detachment onset means operating near a
threshold, where the plasma responds to a small change in the control parameter
by moving a long way. Add the current ramps at the start and end of a discharge,
periodic core relaxations, and deliberate density modulation, and the picture is
of a system that is repeatedly stepped rather than one that is ever allowed to
settle. Whether the intervals between steps are milliseconds or seconds does not
by itself decide the question; what decides it is whether the plasma reaches its
new balance before the next step arrives.

Now set that against the clock the table assumes, and note carefully which clock
that is. It is not the lifetime of an excited state. Those are nanoseconds;
Eq.~\eqref{eq:brightness}'s Einstein coefficients are of order $10^{8}$
\si{\per\second}, and nothing a tokamak does on a millisecond disturbs them.
The relevant clock is the time the \emph{ground-state} population takes to come
back into balance with the ions after conditions change, and that is a much
slower thing.

One can see why it is slower without any model at all. To ionise a hydrogen atom
out of the ground state an electron must bring \SI{13.6}{\electronvolt}. In a
plasma at $\Te = \SI{3}{\electronvolt}$, the fraction of electrons in the
Maxwellian tail with at least that energy is of order
$\exp(-13.6/3) \approx 1.1\times10^{-2}$: about one electron in a hundred. To
excite the same atom to $n=2$ costs only \SI{10.2}{\electronvolt}, and to move an
atom that is \emph{already} in $n=3$ takes a fraction of an electronvolt, which
essentially every electron can supply. The ground state is therefore separated
from everything else by a barrier that only the tail of the distribution can
cross, while the excited levels exchange population freely among themselves. That
is the physical origin of the two very different clocks, and
Chapter~\ref{ch:theory} turns the argument into a calculation.

A chemical engineer will recognise this as a Damk\"ohler-number question, and
that is exactly what it is. If the characteristic time of the disturbance is
short compared with the slowest chemical relaxation in the system, then the
composition during the disturbance is not the equilibrium composition, and a
chart that assumes equilibrium describes a state the system has not reached. Here
the disturbance lasts milliseconds; the fast levels relax in nanoseconds; and the
slow reservoir relaxes on a timescale Chapter~\ref{ch:theory} computes and
Chapter~\ref{ch:results} compares against the event.

So the complication, stated in full: a table built on two assumptions is being
applied to a plasma that plainly violates the second one, and possibly the first.
Which of the two actually fails, and what it costs, is not obvious from anything
said so far, and it turns out not to be obvious at all.

%-----------------------------------------------------------------------------
\section{What is already known}
\label{sec:prior_work}
% QUESTION: who has asked this before, and what did they establish?

\emph{The question: how much of this is new, and how much has been known for
decades?}

This question has been asked before, in several forms, and it is better to be
precise about what was established than to be told at an examination. The
contribution of this thesis is narrower than the question it asks, and the
narrowing is worth stating at the outset.

\paragraph{The time-dependent hydrogen problem has been solved before, and the
overlap is close enough to state plainly.} Sawada and Fujimoto studied exactly
this system: atomic hydrogen, a CR model, and an abrupt change of plasma
parameters. Their paper is titled, in part, for the validity range of the
quasi-steady-state solution of the coupled rate equations
\cite{SawadaFujimoto1994}. What they established is a statement about
\emph{which level controls the approach to steady state}. Solving the coupled
rate equations after a step, they found that in an ionising plasma the levels
below Griem's boundary, the level at which radiative and collisional
depopulation rates are equal, relax at their own natural lifetimes while every
level above it relaxes at the rate of the boundary level itself; in a recombining
plasma the roles invert, the levels above the boundary relaxing on their own
timescales and those below following the boundary. Either way the overall
response of the excited manifold is set by the relaxation time of Griem's
boundary level.

That is a result about the excited levels, and it is the reason this thesis does
not claim to have discovered that the excited manifold responds on a finite
timescale. The difference is in what is being asked, and it is visible in how
the calculation is set up: under ionising conditions they keep $n(1)$ constant
and set the ion density to zero \cite[p.~5566]{SawadaFujimoto1994}. The
reservoir is an input to their problem, not a variable of it. Sawada and
Fujimoto asked when the excited-level solution is valid, and answered it with a
timescale internal to that manifold. The question here is what a \emph{diagnostic
inversion} costs when the ground-state reservoir underneath that solution has not
settled. That is a different quantity, controlled by a different and far slower
clock, and one that can be wrong even when the excited-level solution is
perfectly accurate. Chapter~\ref{ch:results} returns to Griem's boundary
level, which the present model locates independently, and uses it as a check.

The separation of timescales that both papers rest on is older than either, and
Sawada and Fujimoto quote it in its original form: a quasi-equilibrium
population of the excited system is established almost instantaneously
\enquote{without the number densities of free electrons and bare nuclei being
appreciably altered}, because \enquote{the relaxation times for the excited
levels are very much shorter than the relaxation time for the ground levels or
for the free electrons} \cite[p.~5566]{SawadaFujimoto1994}. Chapter~\ref{ch:theory}
measures that ratio as $\Msep$ and Chapter~\ref{ch:results} shows that it does
not predict the diagnostic error, which is the point at which this thesis
departs from the tradition it is quoting.

\paragraph{The general validity question was settled, as a matter of principle,
twenty-five years ago.} Greenland analysed the reduction of CR systems in general
form, in two papers of the same year, and reached two conclusions that bear
directly on this work. The first is that the conditions for a valid reduced CR
model are conditions on the \emph{eigenvectors} of the full matrix of rates, and
that, in his words and in contrast to the conventional view, the
eigenvalues have secondary importance \cite{Greenland2001a}. The second is that
the validity criteria he derives are \emph{not related to the equilibrium
timescales} of the system \cite{Greenland2001b}. The two statements come from
different papers and are cited separately here because they are separately
attributable. Those statements are general, they are prior, and this thesis
does not improve on them as statements. What it does is quantify them for one
specific diagnostic: Chapter~\ref{ch:results} shows that the standard
timescale-separation criterion does not predict the diagnostic error, measures
the error it fails to predict, and maps where in the operating range that error
is large. Turning a general criterion into a number a diagnostician can apply is
the contribution; discovering the criterion is not.

\paragraph{That the ratio depends on the neutral reservoir is current best
practice, not a discovery.} As Section~\ref{sec:the_assumption} noted, Verhaegh
and co-workers already treat the neutral fraction as a free assumed parameter and
propagate its uncertainty by Monte Carlo \cite{Verhaegh2019}. The dependence in
Eq.~\eqref{eq:ratio_three_args} is therefore known. What this thesis adds is a
\emph{closed form} for the sensitivity, an exact expression for how strongly
the ratio responds to a change in the reservoir, derived in
Chapter~\ref{ch:theory}, together with a bound showing that the response can
never exceed a certain size; and a \emph{map} of where in the $(\Te,\nel)$
plane that sensitivity is large enough to matter. A closed form does something a
Monte Carlo cannot: it says which conditions are safe without requiring the
calculation to be repeated.

\paragraph{The classical validity criteria answer a different question.} Griem
asked at what principal quantum number collisions become fast enough to hold a
level in thermal equilibrium with the free electrons, and defined a boundary
level above which the upper part of the manifold is thermalised
\cite{Griem1963, Griem1997}. That is a statement about the internal structure of
the excited levels at one instant. It says nothing about whether the ground state
has settled, which is the question here. Chapter~\ref{ch:results} nonetheless
finds Griem's boundary level turning up where the present model says it should,
which is a useful consistency check rather than a new result.

\paragraph{Real divertor Balmer light is not purely atomic.} One line of work
must be flagged now because it bounds what any conclusion of this thesis can
claim. In detached conditions a substantial part of the observed Balmer emission
does not come from the atomic processes modelled here at all. Hydrogen in a cold
divertor is largely molecular, and molecules open additional routes to excited
atomic levels, among them molecular-assisted recombination (MAR), in which a
molecular ion recombines and dissociates into an excited atom. Analyses of
discharges on the TCV and MAST-U tokamaks attribute a large share of the low-$n$ Balmer light to
these channels near detachment onset \cite{Verhaegh2021, Wijkamp2023}. The model
built in this thesis contains no molecules at all. Chapter~\ref{ch:limits} states
the consequence plainly, and it is the principal reason the results here are
presented as statements about a class of CR model rather than as predictions for
a particular machine.

\medskip

The system has been modelled, the general validity criterion has been stated,
and the dependence on the neutral reservoir is known and handled in practice. What is missing is a quantitative, closed-form answer to the specific
question of how much a two-parameter inversion of the Balmer ratio is in error
when the reservoir has not settled, and a map of where that error lives. That is
the gap this thesis fills.

%-----------------------------------------------------------------------------
\section{The question, and the suspect}
\label{sec:question}
% QUESTION: what exactly is being asked, and what does everybody expect the
% answer to be?

\emph{The question: what is the thesis actually going to compute, and what does
the obvious physical argument predict it will find?}

The question, stated plainly:

\begin{quote}
When a hydrogen plasma at divertor conditions is disturbed on a timescale short
compared with the relaxation of its ionisation balance, how large is the error in
a $(\Te,\nel)$ inversion of the $H_\alpha/H_\beta$ ratio, what physical quantity
sets that error, and where in the operating range is it largest?
\end{quote}

Answering it needs a model that does not assume the answer. This work therefore
builds a time-dependent CR model of atomic hydrogen, every level tracked as a
dynamical variable with orbital angular momentum resolved wherever that
resolution changes the physics, subjects it to an abrupt change in electron
temperature of the kind an ELM delivers, and follows the line ratio in time.
Chapter~\ref{ch:atoms} assembles the atomic data it runs on and
Chapter~\ref{ch:theory} assembles the equations.

\paragraph{What we expected to find.} Section~\ref{sec:the_assumption} identified
two assumptions, and the obvious suspect is the first one. The QSS approximation
is the step that eliminates several dozen dynamical variables by fiat; it is the
step that every review of CR modelling flags; and it is the step whose
chemical-engineering analogue, the Bodenstein approximation, every
undergraduate is warned can fail when the intermediates are not fast enough. In a
plasma whose conditions are reset over a few hundred microseconds by every ELM,
the natural expectation is that some of the excited levels will simply fail to
keep up, the population distribution will lag behind the conditions, and $\Robs$
will be wrong for that reason. That was the expectation this work set out to test, and it is why the
model tracks every level in time rather than assuming any of them are stationary.

\paragraph{What is at stake if that expectation is right.} If the excited levels
lag, then the coefficients $a_p$ and $c_p$ of Eq.~\eqref{eq:two_supplies_intro}
are not functions of $(\Te,\nel)$ during a transient, the whole
population-coefficient framework loses its footing during exactly the events one
most wants to diagnose, and the repair is a genuinely time-dependent inversion.

\paragraph{And if it is wrong?} Then the second assumption is still standing
there, untested, and it is the one the table needs just as much.

Which of the two fails, by how much, and what physical quantity controls it, is
the subject of Chapters~\ref{ch:theory} and~\ref{ch:results}. The answer is not
the one stated above as the expectation; that much can be said here, because
the shape of the argument depends on it. The reason it is not is the reason
this thesis exists.

Two things should be said in the same breath, before any of it. First, the claim
this work defends is conditional: \emph{if} the ground-state neutral population
is frozen for the duration of the disturbance, which requires that neutral
transport and recycling do not renew it faster than the atomic physics does, then
the inversion carries the error computed here. Second, the model is
zero-dimensional, optically thin, and contains no molecules.
Chapter~\ref{ch:limits} examines every one of those boundaries and quantifies
what it costs, rather than listing them as caveats. A result stated without its
scope is the characteristic failure of this subject, and it is not repeated here.

Everything from this point on is construction. The section below says what each
of the remaining chapters contributes to it.

%-----------------------------------------------------------------------------
\section{Outline of the thesis}
\label{sec:outline}
% QUESTION: what does each remaining chapter contribute to the argument?

\emph{The question: where does the argument go from here?}

\textbf{Chapter~\ref{ch:atoms}, \emph{The Atoms}}, establishes the inputs. It
describes everything that can happen to a hydrogen atom in a plasma and how fast
each process runs, sets out the levels the model tracks and the reason orbital
angular momentum is resolved for the lower ones and bundled above them, shows how
a measured collision cross section becomes a rate coefficient, and gives the
provenance and accuracy of every dataset used. It makes a structural argument for
including a process the standard treatments often omit, and is explicit about
where the atomic data are least secure.

\textbf{Chapter~\ref{ch:theory}, \emph{Two Clocks}}, is the conceptual core. It
assembles the CR rate equation from a population balance on each level, checks
the assembly against particle conservation, and then diagonalises the rate matrix
to expose the timescale structure the Boltzmann-factor argument of
Section~\ref{sec:transients} anticipated: one very slow mode, the ground state,
sitting far above a dense band of fast ones. It presents the QSS
approximation as the Bodenstein approximation applied to that fast band, derives
the exact dependence of $\Robs$ on the ground-state reservoir, and defines the two
reference states and the two error measures the rest of the thesis keeps apart.

\textbf{Chapter~\ref{ch:validation}, \emph{Does the Model Work?}}, asks what a
wrong model would fail, and then applies those tests: detailed balance, the
coronal and LTE limits, the timescale hierarchy, and an external benchmark
against independently published population coefficients. It also records the
errors this project's own audit found, what each cost, and how each was caught.

\textbf{Chapter~\ref{ch:results}, \emph{What the Model Says}}, reports the
result: the relaxation in time, the size of each of the two errors, the mechanism
that sets the larger one, a map of where in the $(\Te,\nel)$ plane it lives,
whether it survives averaging over a transient event, and whether the standard
timescale-separation criterion predicts any of it.

\textbf{Chapter~\ref{ch:limits}, \emph{What This Model Cannot Say}}, bounds the
claim: optical thickness, molecular channels, the truncation of the level
manifold, the fixed ion reservoir, and the remaining assumptions, with each
consequence quantified where quantification is possible and stated as open where
it is not.

\textbf{Chapter~\ref{ch:conclusions}} states the criterion in a form another
modeller can apply to their own rate matrix, says where in the divertor operating
range a $(\Te,\nel)$ inversion should be distrusted, and sets out what a
time-dependent inversion would require.

The argument begins with the atoms. Before anything can be said about how fast
the populations change, we need to know what the processes are and how well their
rates are known, which is where Chapter~\ref{ch:atoms} starts.
